\section{Conclusiones}
El análisis realizado permitió identificar que la empresa presenta riesgos importantes relacionados con la integridad, disponibilidad y confidencialidad de la información. Entre los principales problemas detectados se encuentran el uso de sistemas operativos obsoletos, cuentas compartidas, falta de protección eléctrica y ausencia de controles de seguridad más avanzados.

Los riesgos más críticos corresponden a la modificación no autorizada de datos de pesaje, la interrupción total de operaciones por fallas eléctricas y el posible robo de información confidencial de proveedores y facturación. Aunque actualmente existen algunos controles básicos, como respaldos en USB y nube, todavía son insuficientes para garantizar una protección adecuada.

La implementación de políticas de seguridad, controles de acceso, protocolos seguros de comunicación y mecanismos de respaldo permitirá reducir considerablemente las vulnerabilidades detectadas. Asimismo, la inversión en equipos más modernos, licencias originales, antivirus y sistemas de energía protegida contribuirá a mejorar la continuidad.
